\subsection{The literal binary ten-point input and its Klein-four correction}

\begin{result}{A.4.29}{Proved}{The binary sphere quotient has an explicit standard action and geometric fusion correction}
\label{thm:sphere10-binary-fusion}
Let $M=\Gamma_{0,[10]}$ be the discrete orientation-preserving
sphere mapping-class group, let $P=\Gamma_{0,10}$ be its pure
subgroup, and put
\[
 E_2=M/([P,P]P^2),\qquad V=P/([P,P]P^2).
                                                               \tag{BFU.1}
\]
Then $V=\mathbb F_2^{35}$, and the actual quotient sequence is
$1\to V\to E_2\to\mathfrak S_{10}\to1$.
For every odd profinite multiplier $\lambda$ and commutator
proword $f$, the standard ten-strand braid formulas evaluated
in $E_2$ induce the genuine automorphism
\[
 \Theta_\lambda=
 \begin{cases}
  1,&\lambda\equiv1\pmod4,\\
  \jmath,&\lambda\equiv3\pmod4,
 \end{cases}
 \qquad \jmath(\sigma_i)=\sigma_i^{-1}.
                                                               \tag{BFU.2}
\]
In particular this binary action is already fixed by the named
finite generator formulas; no additional pentagon equation is
needed to define it.

More generally, let $\mathcal A$ be the actual subgroup of
$\operatorname{Aut}(E_2)$ inducing the identity both on $V$
and on $\mathfrak S_{10}$.  The exact automorphism map is
\[
 V\times C_2\xrightarrow{\sim}\mathcal A,
       \qquad (w,\epsilon)\longmapsto
                    \operatorname{Inn}(w)\jmath^\epsilon.
                                                               \tag{BFU.3}
\]
Use the ten labels
\[
 A_0=(0,\infty),\quad B_0=(1,-1),\quad C_0=(i,-i),
 \quad D_0=(2,-2,1/2,-1/2),
                                                               \tag{BFU.4}
\]
numbered $0,\ldots,9$ in the displayed order.  The actual
geometric group $H=\langle w\mapsto-w,w\mapsto1/w\rangle$
embeds in $E_2$ by its faithful label permutations.
For every $F\in\mathcal A$ there is a correction in the image
of an actual pure mapping class, followed by either the identity
or the positive half twist interchanging the two labels of $C_0$,
such that the corrected action fixes this $H$ pointwise.
The correction is computed by the explicit linear formula
(BFU.12) below, valid for every section over these same
permutations.  Thus it applies to the specified geometric
section without identifying it with an arbitrarily chosen
solution of the finite complement equations.
\end{result}

\begin{proof}
\par\medskip\noindent\textbf{The actual source quotient.}\quad
Write $e_{ij}$ for the pure braid linking generator on the
unordered pair $\{i,j\}$, with $0\leq i<j\leq9$.  The
pure disk braid group has abelianization freely generated
by these $45$ elements.  One direct derivation successively
forgets the last point of the ordered configuration space.
Its fibre is the plane minus the retained points, with free
fundamental group generated by their meridians.  Local
trivializations move points inside disjoint small discs;
a section places the new point at a positive real coordinate
larger than the sum of all retained absolute values plus one.
The fibration exact sequences split, and the pure braid action
fixes every fibre meridian in abelianization.  Induction gives
the claimed free abelianization and the named generators.

The sphere mapping-class presentation adds to the disk braid
relations the two actual words
\[
 \sigma_1\cdots\sigma_8\sigma_9^2\sigma_8\cdots\sigma_1=1,
 \qquad(\sigma_1\cdots\sigma_9)^{10}=1;
                                                               \tag{BFU.5}
\]
see the checked primary presentation
\cite[Theorem~3.2, printed p.~6]{CA_EndoKamada2016}.
In pure linking coordinates these kill the star of one
vertex and the sum of all edges.  Conjugating the first
word gives every vertex star.  Since these words already
lie in the pure subgroup, their entire normal closures
kill exactly the span of those classes in its abelianization.
It follows, before any change of group, that
\[
 V=\mathbb F_2^{\binom{10}{2}}\big/
 \left\langle
   s_i=\sum_{j\ne i}e_{ij}\ (0\leq i\leq9),\quad
   t=\sum_{i<j}e_{ij}
 \right\rangle.
                                                               \tag{BFU.6}
\]
The star span has dimension nine: a linear combination
with coefficients $a_i$ has edge coefficient $a_i+a_j$,
and is zero precisely when all $a_i$ agree.  The total edge
vector is outside this span, since $a_i+a_j=1$ for all
pairs is incompatible with three distinct vertices.
Thus the relation span has dimension ten and $\dim V=35$.
The quotient map on $V$ and the permutation map give
exactly (BFU.1), rather than only a vector space of the
same dimension.

Here is a complete multiplication model for the extension.
For $\pi\in\mathfrak S_{10}$ let $c_\pi$ be the integral
$0$--$1$ edge vector which is one on $\{i,j\}$, $i<j$,
when $\pi^{-1}(i)>\pi^{-1}(j)$.  It is the crossing vector
of a positive permutation braid.  Represent a braid by its
crossing vector modulo four together with its permutation.
The former has the form $c_\pi+2x$.  Reducing $x$ by
(BFU.6) gives the actual coordinates $(x,\pi)$ of $E_2$,
with multiplication
\[
\begin{gathered}
 (x,\pi)(y,\tau)
   =\bigl(x+\pi y+\kappa(\pi,\tau),\pi\tau\bigr),\\
 \kappa(\pi,\tau)
   =\overline{\frac{c_\pi+\pi c_\tau-c_{\pi\tau}}2}
       \quad\text{in }V.
\end{gathered}                                                  \tag{BFU.7}
\]
The numerator is even, since its parity records only the
final permutation.  Its coboundary identity follows by
cancelling the four integral crossing vectors.  On the pure
subgroup the model is exactly its mod-two abelianization,
so the model has no further kernel.  The two words in
(BFU.5) give precisely the relations (BFU.6).

\par\medskip\noindent\textbf{The finite standard action.}\quad
For $i>1$ the standard second braid argument is
$y_i=\sigma_{i-1}\cdots\sigma_1^2\cdots\sigma_{i-1}$.
Its permutation is the product of a word in adjacent
transpositions with its reverse, hence is the identity.
Both $y_i$ and $\sigma_i^2$ therefore map into $V$.
Their image is abelian, so every commutator proword in
these two arguments has value one.  The standard generator
formula consequently becomes $\sigma_i\mapsto\sigma_i^\lambda$
for every $i$, including the first generator.
The square $\sigma_i^2$ is the nonzero edge class
$\overline e_{i-1,i}$, and has order two.  The generator
$\sigma_i$ thus has order four.  The stated assignment
is the identity or the braid mirror according to
$\lambda\pmod4$.

The mirror is a genuine automorphism: reflection of the
punctured sphere conjugates every chosen positive half twist
to its inverse and preserves the two sphere relations.
It also follows directly by substituting the inverse
generators in the braid relations and (BFU.5).
This proves the actual extension in (BFU.2).  In the
coordinates (BFU.7), negating all crossing numbers gives
\[
 \jmath(x,\pi)=(x+\beta(\pi),\pi),\qquad
 \beta(\pi)=\overline{c_\pi},\qquad
 \beta((i-1\ i))=\overline e_{i-1,i}.
                                                               \tag{BFU.8}
\]
In particular $\jmath$ fixes $V$ pointwise.

\par\medskip\noindent\textbf{All automorphisms with these two identity residues.}\quad
An automorphism fixing $V$ and the permutation quotient has
the unique form $(x,\pi)\mapsto(x+b(\pi),\pi)$.
The multiplication (BFU.7) says exactly that
$b(\pi\tau)=b(\pi)+\pi b(\tau)$; the extension cocycle
cancels on both sides.  Thus this actual automorphism group
is $Z^1(\mathfrak S_{10},V)$, and its pure inner subgroup
is the image of $w\mapsto(\pi\mapsto(1-\pi)w)$.

For completeness, the entire finite cohomology calculation
is retained.  Let $T_i$ be the permutation of edge classes
induced by $(i-1\ i)$ and put $b_i=b((i-1\ i))$.
The Coxeter presentation gives the exact linear equations
\[
\begin{aligned}
 (1+T_i)b_i&=0,\\
 (1+T_j)b_i+(1+T_i)b_j&=0 &&(|i-j|>1),\\
 (1+T_iT_j+T_j)b_i+(T_i+1+T_jT_i)b_j&=0
                                                  &&(j=i+1).
\end{aligned}                                                   \tag{BFU.9}
\]
They are all the presentation relations, so their solution
space is the full cocycle space.  In the ordered edge
quotient (BFU.6), the matrix has $315$ columns and rank
$279$.  The coboundary map from $V$ has rank $35$, and
the explicitly displayed vector
$(\overline e_{01},\overline e_{12},\ldots,
  \overline e_{89})$ is outside that image.  It is precisely
the mirror vector (BFU.8).  These statements are checked
by both packed binary elimination and a separate array
elimination with opposite pivot order.  The complete
matrix, its nullspace basis, the coboundary columns and
the mirror vector are archived.  Hence every cocycle is
uniquely $(1-\pi)w+\epsilon\beta(\pi)$, proving (BFU.3).
The product is direct because $\jmath$ fixes $V$.

\par\medskip\noindent\textbf{The labelled section and its explicit correction.}\quad
Write the two permutations of the geometric group (BFU.4)
as
\[
\begin{aligned}
 a&=(2\ 3)(4\ 5)(6\ 7)(8\ 9),\\
 b&=(0\ 1)(4\ 5)(6\ 8)(7\ 9),\qquad c=(4\ 5).
\end{aligned}                                                   \tag{BFU.10}
\]
The two elements of any section over these permutations
are $(x_a,a)$ and $(x_b,b)$, with $x=(x_a,x_b)\in V^2$.
Their square and commutation equations are
\[
\begin{aligned}
 (1+a)x_a&=\kappa(a,a),\qquad
 (1+b)x_b=\kappa(b,b),\\
 (1+b)x_a+(1+a)x_b&=\kappa(a,b)+\kappa(b,a).
\end{aligned}                                                   \tag{BFU.11}
\]
The actual geometric $H$ supplies one such section, and
its injection follows already from the label action.
The correction below is verified on the full affine
solution set of (BFU.11), so no choice of its origin
affects that geometric section.

Define $B_H:V\to V^2$ by
$B_H(w)=((1+a)w,(1+b)w)$, let $C$ act on both factors
by $c$, and put
$k_c=(\kappa(c,a)+\kappa(a,c),
       \kappa(c,b)+\kappa(b,c))$.
Choose the explicit linear right inverse $R_H$ on
$\operatorname{im}(B_H)$ obtained by binary elimination
with all nonpivot variables zero.  Then the correction is
\[
\begin{gathered}
 \beta_H=(\beta(a),\beta(b)),\qquad
 w(x)=R_H\bigl((1+C)x+C\beta_H+k_c\bigr),\\
 \operatorname{Inn}\bigl(w(x)(0,c)\bigr)\jmath
       \text{ fixes }(x_a,a)\text{ and }(x_b,b).
\end{gathered}                                                  \tag{BFU.12}
\]
Here the two explicit mirror values are
\[
\begin{aligned}
 \beta(a)&=\overline e_{23}+\overline e_{45}
                       +\overline e_{67}+\overline e_{89},\\
 \beta(b)&=\overline e_{01}+\overline e_{45}
       +\overline e_{68}+\overline e_{69}
       +\overline e_{78}+\overline e_{79}.
\end{aligned}                                                   \tag{BFU.13}
\]
The membership needed to apply $R_H$ is verified on a
particular solution of (BFU.11) and on every vector in
its full $30$-dimensional homogeneous solution space.
The equation matrix has rank $40$, while $B_H$ has rank
$23$.  Thus the quotient of sections by pure conjugacy
has $2^7=128$ elements.  The computed action of the
half twist $(0,c)$ is translation by the same nonzero
class $[\beta_H]$ on every one of these $128$ classes.
The program separately verifies (BFU.12) using the actual
extension multiplication, for the affine origin and all
thirty homogeneous directions.  Affine linearity then
proves it for every section.

One can also see the formula directly once that linear
membership is known.  Conjugation by $(0,c)$ sends
$(y,a)$ to
$(cy+\kappa(c,a)+\kappa(a,c),a)$, since $c$ commutes
with $a$; conjugation by $w\in V$ adds $(1+a)w$.
Inserting $y=x_a+\beta(a)$ gives exactly the first
component of (BFU.12), and the calculation for $b$ is
identical.  Every vector $w$ is the image of the literal
pure braid obtained by multiplying its chosen edge
generators $A_{ij}$.  The element $(0,c)$ is the image
of the positive half twist $\sigma_5$, supported in a
disc containing only the two labels of $C_0$.  Forgetting
those two labels kills this half twist.  Thus the
correction in (BFU.12) is supplied by actual geometric
source words, not an arbitrary automorphism of the
finite module.

Finally decompose any $F\in\mathcal A$ using (BFU.3).
A pure correction removes its $\operatorname{Inn}(w)$
factor; when its remaining factor is $\jmath$, use
(BFU.12).  This proves the asserted pointwise fusion.
The same calculation records the actual finite
centralizer sequence in the preimage of the colored
permutation centralizer $C_2^5$:
\[
 1\longrightarrow\mathbb F_2^{12}
   \longrightarrow C_{E_2}(H)
   \longrightarrow
   \{(r_A,r_B,r_C,d_a,d_b):r_A+r_B+r_C=0\}
   \longrightarrow1.
                                                               \tag{BFU.14}
\]
Indeed the kernel is $V^H$, of dimension twelve, and
the same $128$-class calculation gives the displayed
stabilizer for every section.  No equality of this
finite centralizer with the image of the full discrete
centralizer is used here.

The program
\path{work/ghost_code/check_sphere10_binary_extension.py}
retains the quotient basis, all colored cocycle entries,
the full Coxeter equation matrix, all section classes,
all restriction maps, and the correction formula in
every affine direction.  Its independent array ranks,
the explicit mirror cocycle and the actual pure-braid
correction words are asserted in the executed output.
The binary standard action (BFU.2) is now an established
finite input.  Lifting it to the relative ternary cover,
and connecting that larger input to the required finite
pentagon comparison, remain additional tasks; no such
lift is asserted by the binary result.
\end{proof}
